% SPDX-FileCopyrightText: Copyright (c) 2026 NVIDIA CORPORATION & AFFILIATES. All rights reserved.
% SPDX-License-Identifier: Apache-2.0
%
% ROLE: why worker selection is a deployment decision worth learning, why it is hard to settle by
%   hand or by experiments on GPUs, what the campaign does about it, how "better" is decided, and
%   what was found, with its scope.
% EVIDENCE LEVELS: latencies quoted here are AIS performance-model timings, not GPU measurements.
%   Headline numbers: CR/facts/test_results.json; scope: report_must_state.
% STATUS: final, including the live GPU check (sec:live).

\section{Introduction}
\label{sec:intro}

\paragraph{The decision.}
A KV-aware router places each request on the worker whose cache already holds the longest prefix
of its prompt, unless that worker is busy enough that recomputing the prefix elsewhere is cheaper.
On the deployment studied here, Qwen3-32B on vLLM with two H100 GPUs per worker, both sides of this
trade are large in \ais{}'s timing model. Reuse saves prefill: a \num{120000}-token prompt has a
time to first token of \SI{17.8}{s} on an idle worker that lacks its prefix and \SI{56}{ms} on an
idle worker that already holds the first \num{119808} tokens.
% src: CR/audits/setup/determinism-cost-context-r1.md "New: long prefix reuse at 120K"
%      (17,815.54 ms without the prefix; TTFT 56.06 ms with reused_input_tokens 119,808);
%      CR/facts/setup.json long_context.evidence.single_request_ttft_ms_ais."120000"
Load costs decode speed: one decode step at \num{8192} tokens of context takes \SI{15.3}{ms} at
batch size 1 and \SI{109.9}{ms} at batch size 256, so each request placed on a worker slows the
token stream of every other request there.
% src: CR/facts/setup.json ais_active_evidence.decode_step_ms_by_batch_and_context
%      (b1_ctx8192 15.3367 ms, b256_ctx8192 109.935 ms)
Dynamo's default router makes the trade with a fixed cost, the prefill blocks a worker would still
have to compute plus its decode load (\cref{eq:default-cost}). Published routers draw the line
elsewhere. LMetric multiplies a worker's prefill work by its number of running requests instead of
adding the two \citep{zhang2026lmetric}. DualMap hashes a prefix of each prompt to two candidate
workers and picks between them by their current state \citep{yuan2026dualmap}. Lodestar narrows its
candidates to two consistent-hash choices in the same way, but only while cluster memory
utilization is above a threshold, 80\% in its evaluation \citep{lim2026lodestar}. SMetric keeps a
session's follow-up turns on the worker that holds its context while that worker can still meet a
latency deadline \citep{wang2026smetric}.
% src: CR/literature/LESSONS.md LR-04 (LMetric P-token x batch; SMetric Fig 13 rule) and LR-07
%      (Lodestar k = 2 filter, pdf p.6; DualMap pdf p.5-6), all VERIFIED

These rules disagree, and which one wins depends on the traffic. LMetric's authors report that the
product beats linear scores tuned for each workload, Dynamo's among them, on their traces
\citep{zhang2026lmetric}, whose ideal prefix reuse SMetric's authors put at about 50--70\%. On an
agentic workload with 82\% ideal reuse and no global KV tier, which is the setting of our replay,
SMetric's authors measured LMetric 36\% below a tuned linear score in throughput within SLO,
because its reuse ratio fell to 45\% against 74\% \citep{wang2026smetric}. Even the best linear
weight moves: LMetric's own sweep put it at 0.7 on one workload and 0.55 on another with a similar
access pattern \citep{zhang2026lmetric}.
% src: CR/literature/LESSONS.md LR-07 evidence (SMetric pdf p.6; LMetric pdf p.6, p.11), VERIFIED
An operator has to choose one rule and tune it for a deployment whose traffic may change.

\paragraph{Why the choice is hard to settle.}
Three properties of the problem make the comparison expensive.
\begin{enumerate}
  \item \emph{No per-decision label.} A routing decision changes the cache contents and queues
    that later requests meet, so its quality shows only in aggregate, over minutes of traffic at a
    realistic load. A policy can be scored only on whole episodes.
  \item \emph{Small differences next to workload variation.} On this deployment, the two strongest
    tuned heuristics finished within 0.002 of each other in mean validation score, and their order
    flipped between the mean over cells ($+0.0018$) and the mean over independent segments
    ($-0.0087$) (\cref{sec:res-selection}). Detecting effects of a few percent needs many cells,
    replicates and independent workload segments.
    % src: CR/facts/finalists.json rankings_val_k3_10.baseline_pool (ramjet 0.14113, llmdpp 0.13934);
    %      CR/facts/tierBC.json decisions.val_best_baseline ("margin +0.0018 cell mean (SE 0.0061,
    %      9/14 cells) but -0.0087 segment mean (3/6 segments): a near tie")
  \item \emph{Regime dependence.} In classical load balancing, how much waiting a dispatch
    policy can avoid depends on the number of servers and on the load
    \citep{vanderboor2022survey}. In a measured study of learned routing for disaggregated
    serving, the learned router tied join-the-shortest-queue at decode width 3, led at width 4
    and tied round-robin again at width 6 \citep{tumkur2026calibrate}. A policy tuned at
    one worker count can therefore lose at another.
    % src: CR/literature/LESSONS.md LR-12 evidence (CtR pdf p.4-5), VERIFIED
\end{enumerate}
Settling the choice by measurement means many long runs. LMetric's authors note that they could
not sweep all configurations because replaying each trace on their testbed consumes substantial GPU
time \citep{zhang2026lmetric}. The tuning in this study is larger: 66 CMA-ES restarts of 400
evaluations each, every evaluation replaying all 34 train cells on 2 workload replicates, plus
validation. Each train cell simulates 6 to 29 minutes of serving on 4 or 8 workers, that is, 8 to
16 GPUs.
% src: CR/audits/final/completeness.md and refute-headline-2.md (66 restarts, 400 evals, 34 train
%      cells, K = 2); CR/cells/train.jsonl measure.warmup_ms + measure.window_ms over the
%      arrival-basis cells (6.3 to 29.0 simulated minutes); N in {4, 8} at TP2.
% src: CR/literature/txt/llm-routing/zhang2026-lmetric.txt lines 767-772 ("we cannot afford to sweep
%      all possible configurations, as replaying each trace on the testbed consumes a substantial
%      amount of GPU time")

\paragraph{Approach.}
We first describe the router's worker-selection API (\cref{sec:background}). A routing policy is a
YAML instance of a catalog type that composes filters, scorers and a picker over inputs it declares:
cache overlap, load, capacity and request facts. Dynamo's default cost function and every ported
heuristic in this study are such instances, and so is the learned policy, so all of them are
configured, tuned and replayed through the same interface. We then use \aisim{} offline replay as a
development environment (\cref{sec:simulator}). It runs Dynamo's router, including the builtin policy
catalog the live router loads, against mock vLLM engines in simulated time, and times every forward
pass with AISimulate's performance model for the configured model, GPU and parallelism. A
\num{2000}-request trace replays in about \SI{1.5}{s} of one CPU core, and identical inputs replay
identically, so competing policies can be compared on common random numbers.
% src: CR/facts/setup.json replay_cost (about 1.5 s per 2,000-request replay in a fresh process);
%      CR/CONTRACT.md A1 (identical re-runs are deterministic)
In that environment we tune a worker-selection function, \lc{} (\cref{sec:model}). It is a
conditional-logit choice model: it scores every candidate worker with one shared utility, linear
in features the live router already computes, and picks the highest. CMA-ES tunes the
coefficients on windowed SLO goodput (\cref{sec:training}). The simulator only trains and scores:
the headline policy reads no performance-model signal, so the YAML file that training writes runs
in the live router unchanged. A separately reported league of variants may read AISimulate
estimates at runtime (\cref{sec:res-ais}).
% src: CR/CONTRACT.md A16 (AIS league, secondary); WT/lib/router-plugins/builtin/src/learned_choice/FEATURES.md

The choice-model form gives three properties the study relies on. It is defined for any candidate
set, so a policy tuned at $\Nw \in \{4, 8\}$ can run at $\Nw = 2$ or 32 and be tested there. It
contains the default router: the coefficient vector $\theta = -\basis{0}$ reproduces the default
router's choice (\cref{sec:parity}), so training starts from the status quo and every learned
coefficient reads as a change to it. And it is cheap: a decision costs $O(\Nw\dfeat)$ operations
for $\dfeat$ features (8 in the base set, 23 in the extended set that won).

\begin{figure}[tbp]
  \centering
  \begin{tikzpicture}[
      node distance=6mm and 7mm,
      box/.style={draw, rounded corners=2pt, align=center, font=\small, inner sep=4pt, minimum height=9mm},
      arr/.style={-{Stealth[length=2mm]}, thick}]
    \node[box] (cells) {frozen cells\\\scriptsize trace, transform, load, $\Nw$, SLO};
    \node[box, right=of cells] (replay) {\aisim{} replay\\\scriptsize router + mock vLLM, \ais{} timing};
    \node[box, right=of replay] (score) {scorer\\\scriptsize $\good_j$, $\Gw_{\cell,\krep}$, $\clr$};
    \node[box, below=of score] (cma) {CMA-ES\\\scriptsize $\Jobj(\pol)$, CRN replicates};
    \node[box, below=of replay] (yaml) {policy YAML\\\scriptsize \lc{} coefficients};
    \node[box, below=of cells, fill=black!5] (live) {live KV router\\\scriptsize same YAML, no \ais{}};
    \draw[arr] (cells) -- (replay);
    \draw[arr] (replay) -- (score);
    \draw[arr] (score) -- (cma);
    \draw[arr] (cma) -- (yaml);
    \draw[arr] (yaml) -- (replay);
    \draw[arr, dashed] (yaml) -- node[above, font=\scriptsize] {ship} (live);
  \end{tikzpicture}
  \caption{The training loop. A policy is a YAML file read by the router's builtin policy catalog,
    so each CMA-ES candidate costs a new YAML and a set of replays, with no rebuild. \ais{} times
    the simulated engines and so shapes the score; the headline policy never reads it.}
  \label{fig:loop}
\end{figure}

\paragraph{How ``better'' is decided.}
The comparison is on equal footing. Every policy, the learned arms and each of 11 tuned heuristic
baselines, is one configuration tuned on the pooled train cells with the same CMA-ES budget and
restarts and selected on validation; the best baseline, and by the same rule the headline learned
arm, are chosen on fresh validation replicates (\cref{sec:baselines,sec:eval}).
% src: CR/CONTRACT.md A2.1, A17.2; CR/facts/finalists.json selection
The comparison is made once, on a test set frozen with SHA-256 hashes before any tuning. Its 60
cells come from 12 independent workload segments and hold out time windows, agent plays, session
seeds, a trace family, transform ranges and worker counts. The headline is a one-sided exact
Wilcoxon signed-rank test over the segments, registered before training. Reported but not tested
are robustness to the SLO scale and definition, to stale router state (replay gives the router
perfectly fresh cache and load information, a live router does not) and to engine timing, and a
check of the ranking on GPUs.
% src: CR/facts/HEADLINE_TEST.json; CR/facts/test_freeze.json; CR/CONTRACT.md A5, A13, A14

\paragraph{What we found.}
% src: CR/facts/test_results.json headline, report_must_state; CR/facts/fairness_in_class_heuristic.json;
%      CR/facts/robustness.json; CR/audits/phase2-mid/gaming-concentration.md
Stated as the final audit allows: in AIS-timed simulation, the learned router M1-v2 beats the
branch's ported heuristics, each tuned with an equal budget. Against the validation-selected best
baseline, tuned \policy{ramjet}, it is ahead on 10 of the 12 test segments and on all 8 informative
ones (the 4 session segments tie at ceiling), with a segment-mean clipped log-ratio of $+0.0316$
and one-sided exact Wilcoxon $p = 0.0017$ (\cref{sec:res-headline}). It keeps first place among the
robustness finalists under router-state lag of 10 to \SI{200}{ms} and under four engine-timing
perturbations, and it ranks first under all four alternative definitions of a good request. Four
qualifications travel with the result.
\begin{itemize}
  \item Only the sign is a claim: the effect is below the pre-registered minimum detectable effect
    (0.038) and below the spread the timing perturbations induce (0.060).
  \item It is not ``learning beats every heuristic''. The baseline pool holds the branch's ports as
    implemented, and its \policy{lmetric} port lacks LMetric's queued-prefill term (a fix made to the
    port after the campaign froze adds it; \cref{sec:ports}). Faithful LMetric
    lies inside the learned model's class; untuned, it beats every tuned baseline on validation, and
    with the default cost it reproduces 53--66\% of the learned model's validation margin
    (\cref{sec:res-fairness}).
  \item It is a request-count property. A token-weighted goodput is mixed, and prompts of 32K
    to 64K tokens see higher time to first token than under \policy{ramjet}
    (\cref{sec:gaming}).
    % src: CR/facts/test_results.json report_must_state MS-F1-ttft (32K-64K p90 +0.135, 11/12; >= 64K
    %      filtered +0.021, 3/7, p 1.0, so the TTFT statement stops at 64K)
  \item Black-box search games a request-count objective. The first tuned model kept a worker
    whose requests were all bad; constraining load coefficients so that load cannot attract
    removed it, and every later learned arm runs under that constraint (\cref{sec:gaming}).
\end{itemize}
A pre-registered live check then served six of the $\Nw = 4$ test cells on H100 GPUs with Dynamo's
router running the same policy files (\cref{sec:live}). The live ranking of the six policies checked
agrees with simulation (pooled Kendall $\kendall$-b 1.00), and M1-v2 leads \policy{ramjet} in each of the
five cells where simulation predicts a lead, each time by more than the registered noise bar. The
check validates the ranking and that sign only: with one run per policy on six cells chosen after the
test pass, it says nothing about the headline's magnitude, and live goodput departs from simulation
by policy-dependent factors of 0.72 to 2.18, enough that on one cell the simulated gain over the
default router is not reproduced beyond live noise.
% src: CR/facts/live_results.json wording, primary.kendall_tau.k0 (tau_b 1.0), primary.sign_agreement
%      (5 of 5, beyond_noise_agree 5); CR/report/LIVE.md 15.2 (0.72-2.18; sessions not reproduced beyond
%      live noise), 15.7

\paragraph{Contributions.}
\begin{itemize}
  \item A description of Dynamo's worker-selection API as routing primitives, with the default
    cost function and nine other catalog policy types written in that vocabulary, and a list of the
    major caveats the campaign uncovered in the API, the replay host and the workloads, each with
    its evidence and status (\cref{sec:background}).
  \item \lc{}, a deployable worker-selection policy for Dynamo's router plugin catalog: a
    discrete-choice model over router-observable features that contains the default router as one
    coefficient setting and reads no simulator signal (\cref{sec:model}).
  \item A tuning method for routing policies in replay: CMA-ES on an episodic, windowed goodput
    objective with common-random-number workload replicates, applied with the same budget,
    restarts and selection rule to the learned model and to every heuristic baseline
    (\cref{sec:training,sec:baselines}).
  \item An objective and evaluation protocol for routing studies: goodput defined by inter-token
    latency and end-to-end slowdown over a request's own uncontended latency, measured in a window
    that excludes warm-up and drain; held-out worker counts; a pre-registered test over independent
    workload segments; and staged adversarial audits (\cref{sec:setup,sec:eval}).
  \item Results with their scope: the headline and its strata, worker-count extrapolation, the
    model ladder, the learned coefficients, robustness, the objective-gaming finding and its
    remedy, the secondary metrics that bound the claim, and a pre-registered live GPU check of the
    ranking (\cref{sec:results,sec:gaming,sec:live}).
\end{itemize}

\paragraph{Reading guide.}
\Cref{sec:background} introduces the routing primitives, the policies and the caveats.
\Cref{sec:setup} defines the decision and the objective, \cref{sec:simulator} the simulator and
workloads, \cref{sec:model} the policy and \cref{sec:training} its training. \Cref{sec:baselines}
describes how the baselines are tuned, \cref{sec:eval} the evaluation protocol, \cref{sec:results}
the results and \cref{sec:gaming} the concentration findings. \Cref{sec:live} reports the live GPU
validation. \Cref{sec:threats} discusses limitations and threats to validity, and
\cref{sec:related} relates the work to prior research. Labels of the form A$n$ refer to the
campaign contract's numbered amendments and LR-$n$ to the numbered lessons of its literature
review; both are part of the campaign record (\cref{app:repro}). Every number traces to a campaign
record named in a comment in the \LaTeX{} source.
